\slcol{श्रीभगवानुवाच —\\
\Index{मय्यासक्तमनाः} पार्थ योगं युञ्जन्मदाश्रयः । \\
असंशयं समग्रं मां यथा ज्ञास्यसि तच्छृणु ॥ ७.१ ॥}
\translate{śrībhagavānuvāca —\\
mayyāsaktamanāḥ pārtha yōgaṁ yuñjanmadāśrayaḥ | \\
asaṁśayaṁ samagraṁ māṁ yathā jñāsyasi tacchr̥ṇu || 7.1 || }
\cquote{Sri Bhagavan said: 
O Arjuna, centre your mind upon Me, seek refuge in Me, and engage in the discipline of \emph{yoga}. In this way, you shall come to know Me completely, without any doubt. Now, hear how this knowledge unfolds.}

\slcol{\Index{ज्ञानं तेऽहं सविज्ञानमिदं} वक्ष्याम्यशेषतः । \\
यज्ज्ञात्वा नेह भूयोऽन्यज्ज्ञातव्यमवशिष्यते ॥ ७.२ ॥ }
\translate{jñānaṁ tē:'haṁ savijñānamidaṁ vakṣyāmyaśēṣataḥ | \\
yajjñātvā nēha bhūyō:'nyajjñātavyamavaśiṣyatē || 7.2 || }
\cquote{I shall now impart to you complete knowledge together with practical wisdom (realisation). When you grasp this fully, nothing else in this world will remain for you to know.}

\newpage
\begin{mananam}{\mananamfont Mananam Śloka - 2}
\normalfont\mananamtext Am I inclined to seek spiritual knowledge without applying it in practice? As in other areas of life, what distinguishes theoretical knowledge from practical understanding in spiritual life? How can my spiritual insights translate into wisdom and personal realisation, making them applicable in everyday life?
\end{mananam}
\sreflect
\begin{inspiration}{\mananamfont\normalfont\textbf{ Inspiration}}
\normalfont\mananamtext In every aspect of life, there is a distinction between theoretical and practical knowledge. While spiritual knowledge can have general applications, the way each individual applies it in his or her own life will vary. However, there is a key difference in spiritual knowledge: unlike other fields where the accumulation of knowledge is limitless, spiritual knowledge has an endpoint. \emph{Brahma Vidyā} is described as that knowledge by which all other knowledge becomes known—knowing One leads to knowing all!
\end{inspiration}
\newpage

\slcol{\Index{मनुष्याणां सहस्रेषु} कश्चिद्यतति सिद्धये । \\
यततामपि सिद्धानां कश्चिन्मां वेत्ति तत्त्वतः ॥ ७.३ ॥ }
\translate{manuṣyāṇāṁ sahasrēṣu kaścidyatati siddhayē | \\
yatatāmapi siddhānāṁ kaścinmāṁ vētti tattvataḥ || 7.3 || }
\cquote{Among thousands of people, only a few strive for perfection. And even among those rare ones who reach perfection, only one in thousands truly perceives My essence.}

\slcol{\Index{भूमिरापोऽनलो वायुः} खं मनो बुद्धिरेव च । \\
अहङ्कार इतीयं मे भिन्ना प्रकृतिरष्टधा ॥ ७.४ ॥ }
\translate{bhūmirāpō:'nalō vāyuḥ khaṁ manō buddhirēva ca | \\
ahaṅkāra itīyaṁ mē bhinnā prakr̥tiraṣṭadhā || 7.4 || }
\cquote{Earth, water, fire, air, space, mind, intellect, and ego—these eight constitute My separated material energies.}

\slcol{\Index{अपरेयमितस्त्वन्यां} प्रकृतिं विद्धि मे पराम् । \\
जीवभूतां महाबाहो ययेदं धार्यते जगत् ॥ ७.५ ॥ }
\translate{aparēyamitastvanyāṁ prakr̥tiṁ viddhi mē parām | \\
jīvabhūtāṁ mahābāhō yayēdaṁ dhāryatē jagat || 7.5 || }
\cquote{O mighty-armed Arjuna, beyond this lower \emph{prakṛti}, know that there exists a higher \emph{prakṛti} of Mine. It is the source of all living beings, and through it the entire universe is upheld.}

\slcol{\Index{एतद्योनीनि भूतानि} सर्वाणीत्युपधारय । \\
अहं कृत्स्नस्य जगतः प्रभवः प्रलयस्तथा ॥ ७.६ ॥ }
\translate{ētadyōnīni bhūtāni sarvāṇītyupadhāraya | \\
ahaṁ kr̥tsnasya jagataḥ prabhavaḥ pralayastathā || 7.6 || }
\cquote{Know that all beings arise from these two \emph{prakṛtis’} wombs. I am the origin of the entire universe, and I am also its dissolution.}

\newpage
\begin{mananam}{\mananamfont Mananam Śloka - 3}
\normalfont\mananamtext Can I observe in my own society and surroundings that many people are focused on pursuing material goals? Can I see that others are primarily engaged in achieving personal success in their respective fields? Why is it that so few pursue genuine self-development and spiritual growth? Even among those who seek to discover their true Self or connect with the Divine, why do so few manage to persevere to the end? In my own life, what are the obstacles preventing me from choosing the higher path to lasting happiness?
\end{mananam}
\sreflect
\begin{inspiration}{\mananamfont Inspiration}
  \normalfont\mananamtext One of the \emph{Upaniṣads} teaches that all human beings have a choice in their life’s path. Some choose the path that pleases the senses and feeds the ego—this is the path of \emph{preyas}. Others commit to the path of the higher good, which leads to ultimate happiness—this is the path of \emph{śreyas}. Although many people are aware of these paths, very few choose to walk the latter. And even among those who do, only a few persist. This is where the guidance of the scriptures, \emph{gurus}, and the support of the right companions become essential for a seeker (\emph{sādhaka}). They are your most loyal cheerleaders!
\end{inspiration}
\newpage

\begin{mananam}{\mananamfont {Mananam Śloka - 4, 5}}
\normalfont\mananamtext
What are the various forces in nature that we experience? From a scientific perspective, we categorise them as matter, energy, and even mind. Are these forces in nature independent, or are they interdependent?
Could there be another force distinct from these, yet sustaining them both? Can I directly observe the workings of these different forces in nature without relying on concepts from science or religion?
\end{mananam}
\sreflect
\begin{inspiration}{\mananamfont Inspiration}
\normalfont\mananamtext The \emph{Gita} is pointing to us that \emph{para} and \emph{apara} \emph{prakṛti} are the two forces at work in nature. While the functions of mind and matter are evident to all, the role of Consciousness, which supports and holds all matter together, remains a profound mystery that humans are meant to unravel. As long as we are content with a life identified with matter, or at most, with the mind, our true identity will remain hidden from us. A \emph{jijñāsu} is someone who sincerely desires to know this deeper reality, while a \emph{mumukṣu} is one who is determined to pursue it at any cost.
\end{inspiration}
\newpage

\slcol{\Index{मत्तः परतरं} नान्यत्किञ्चिदस्ति धनञ्जय । \\
मयि सर्वमिदं प्रोतं सूत्रे मणिगणा इव ॥ ७.७ ॥ }
\translate{mattaḥ parataraṁ nānyatkiñcidasti dhanañjaya | \\
mayi sarvamidaṁ prōtaṁ sūtrē maṇigaṇā iva || 7.7 ||}
\cquote{O Dhananjaya (the winner of wealth, Arjuna), there is nothing higher than Me. All that exists in the universe is woven into Me, as pearls are strung upon a thread.}

\slcol{\Index{रसोऽहमप्सु कौन्तेय} प्रभास्मि शशिसूर्ययोः । \\
प्रणवः सर्ववेदेषु शब्दः खे पौरुषं नृषु ॥ ७.८ ॥ }
\translate{rasō:'hamapsu kauntēya prabhāsmi śaśisūryayōḥ | \\
praṇavaḥ sarvavēdēṣu śabdaḥ khē pauruṣaṁ nr̥ṣu || 7.8 || }
\cquote{O Kaunteya, I am the very essence of water, the radiance of the sun and the moon, the sacred syllable \emph{Om} in all the Vedas, the sound that pervades space, and the essential strength found in human beings.}


\slcol{\Index{पुण्यो गन्धः पृथिव्यां} च 
तेजश्चास्मि विभावसौ । \\
जीवनं सर्वभूतेषु 
तपश्चास्मि तपस्विषु ॥ ७.९ ॥ }
\translate{puṇyō gandhaḥ pr̥thivyāṁ ca 
tējaścāsmi vibhāvasau | \\
jīvanaṁ sarvabhūtēṣu 
tapaścāsmi tapasviṣu || 7.9 || }
\cquote{I am the pure fragrance of the earth, the brilliance within fire, the life force that dwells in all beings, and the austerity of those who practise penance.}

\slcol{\Index{बीजं मां सर्वभूतानां} विद्धि पार्थ सनातनम् । \\
बुद्धिर्बुद्धिमतामस्मि तेजस्तेजस्विनामहम् ॥ ७.१० ॥ }
\translate{bījaṁ māṁ sarvabhūtānāṁ viddhi pārtha sanātanam | \\
buddhirbuddhimatāmasmi tējastējasvināmaham || 7.10 || }
\cquote{O Partha, know that I am the eternal seed of all beings. I am the wisdom that shines in the wise, and the radiance that illumines the radiant.}

\newpage
\begin{mananam}{\mananamfont Mananam Śloka - 7}
\normalfont\mananamtext What do I understand to be the supreme force that governs all other forces of nature? How do I conceive of this essential energy that sustains everything, serving as the foundation of all creation? What insights do the scriptures provide about this? Where and how can I see, know, or comprehend it? How can I perceive the workings of this central force within my own body and mind? Which spiritual practices might help me directly experience it?
\end{mananam}
\sreflect
\begin{inspiration}{\mananamfont Inspiration}
\normalfont\mananamtext A necklace of pearls is held together by a thread, yet our attention is drawn to the beads, not the thread itself. Similarly, the Lord is the unseen force, the source of all other forces at work in the universe. For the spiritually aware, this Divine Energy is the essence of everything, though unseen—just as gravity is constantly at work, yet never directly perceived.
\end{inspiration}
\newpage
\begin{mananam}{\mananamfont Mananam Śloka - 10}
\normalfont\mananamtext
What is it within the tiny seed that allows it to grow into a giant tree? What gives the brain its ability to generate thought and expression? What illuminates the sunlight and the light of the stars? What infuses beauty into some of the most magnificent wonders of nature? Have I lost the curiosity and wonder of a child? Is there value in looking at nature with awe? Could such admiration lead me to discover the One Divine hiding behind it all?
\end{mananam}
\sreflect
\begin{inspiration}{\mananamfont Inspiration}
\normalfont\mananamtext The beauty and complexity of creation cannot be a mere accident. The force that holds the planets and stars in place, that governs our digestion and respiration, must be the most noble presence in all of existence. If this presence were to demand payment for its services, we might be enslaved for eternity. Yet, the Lord serves us without expecting anything in return. Can we not at least express our gratitude for such unconditional love?
\end{inspiration}
\newpage

\slcol{\Index{बलं बलवतां चाहं} कामरागविवर्जितम् । \\
धर्माविरुद्धो भूतेषु कामोऽस्मि भरतर्षभ ॥ ७.११ ॥ }
\translate{balaṁ balavatāṁ cāhaṁ kāmarāgavivarjitam | \\
dharmāviruddhō bhūtēṣu kāmō:'smi bharatarṣabha || 7.11 || }
\cquote{I am the strength of the strong, free from desire and attachment. And I am the desire that dwells in all beings, when it is in harmony with \emph{dharma}, O \emph{Bharatarṣabha} (the best of the Bharata family).}

\slcol{\Index{ये चैव सात्त्विका} भावा राजसास्तमसाश्च ये । \\
मत्त एवेति तान्विद्धि न त्वहं तेषु ते मयि ॥ ७.१२ ॥ }
\translate{yē caiva sāttvikā bhāvā rājasāstamasāśca yē | \\
matta ēvēti tānviddhi na tvahaṁ tēṣu tē mayi || 7.12 || }
\cquote{All beings and objects that embody the qualities of \emph{sattva} (goodness), \emph{rajas} (passion), and \emph{tamas} (ignorance)—know that these arise from Me. Yet I am not contained within them; rather, they abide in Me.}

\slcol{\Index{त्रिभिर्गुणमयैर्भावैरेभिः} सर्वमिदं जगत् । \\
मोहितं नाभिजानाति मामेभ्यः परमव्ययम् ॥ ७.१३ ॥ }
\translate{tribhirguṇamayairbhāvairēbhiḥ sarvamidaṁ jagat | \\
mōhitaṁ nābhijānāti māmēbhyaḥ paramavyayam || 7.13 || }
\cquote{The entire world is deluded by the states of three \emph{guṇas—sattva}, \emph{rajas}, and \emph{tamas}, and thereby do not realise Me, the imperishable reality, who stands beyond these qualities.}

\slcol{\Index{दैवी ह्येषा गुणमयी} मम माया दुरत्यया । \\
मामेव ये प्रपद्यन्ते मायामेतां तरन्ति ते ॥ ७.१४ ॥ }
\translate{daivī hyēṣā guṇamayī mama māyā duratyayā | \\
māmēva yē prapadyantē māyāmētāṁ taranti tē || 7.14 || }
\cquote{Truly, this divine \emph{māyā} of Mine, woven from the three \emph{guṇas}, is exceedingly difficult to overcome. Yet those who surrender unto Me cross beyond this \emph{māyā}.}

\newpage
\begin{mananam}{\mananamfont Mananam Śloka - 14}
\normalfont\mananamtext What is the nature of delusion? Could it be that we are all caught in a mass illusion, as the Hindu scriptures suggest? What are some of the actions that I have done in my life while in a deluded state of mind, only to regret later? What helps me return to a state of awareness? What habits do I still indulge in that may stem from a deluded mind? How can I avoid falling into such traps again? Is there any external guidance or support that can help me break free from these states?
\end{mananam}
\sreflect
\begin{inspiration}{\mananamfont Inspiration:}
\normalfont\mananamtext \emph{Māyā} is one of the most profound concepts across world philosophies, describing the force that disconnects us from our true nature and the Divine. Modern examples like movies and virtual reality provide a tangible way to grasp the idea of \emph{māyā}. Just as a magician manipulates and controls a magic show, there is a divine magician who governs the illusive power of \emph{māyā}. The purpose of the magician is not based on some cruel or manipulative intent. Those who seek refuge in Him are lifted into reality, where the illusions and nightmares end forever.
\end{inspiration}
\newpage

\slcol{\Index{न मां दुष्कृतिनो} मूढाः प्रपद्यन्ते नराधमाः । \\
माययापहृतज्ञाना आसुरं भावमाश्रिताः ॥ ७.१५ ॥ }
\translate{na māṁ duṣkr̥tinō mūḍhāḥ prapadyantē narādhamāḥ | \\
māyayāpahr̥tajñānā āsuraṁ bhāvamāśritāḥ || 7.15 || }
\cquote{The wicked, the deluded, and the lowest among men do not turn to Me. Those whose understanding is clouded by illusion (\emph{māyā}), and who embrace demonic ways, likewise do not surrender to Me.}

\slcol{\Index{चतुर्विधा भजन्ते मां} जनाः सुकृतिनोऽर्जुन । \\
आर्तो जिज्ञासुरर्थार्थी ज्ञानी च भरतर्षभ ॥ ७.१६ ॥ }
\translate{caturvidhā bhajantē māṁ janāḥ sukr̥tinō:'rjuna | \\
ārtō jijñāsurarthārthī jñānī ca bharatarṣabha || 7.16 || }
\cquote{O \emph{Bharatarṣabha} (Arjuna), four kinds of virtuous ones worship Me: those who are distressed, those seeking knowledge, those desiring worldly gains, and those endowed with wisdom.}

\slcol{\Index{तेषां ज्ञानी नित्ययुक्त} एकभक्तिर्विशिष्यते । \\
प्रियो हि ज्ञानिनोऽत्यर्थमहं स च मम प्रियः ॥ ७.१७ ॥ }
\translate{tēṣāṁ jñānī nityayukta ēkabhaktirviśiṣyatē | \\
priyō hi jñāninō:'tyarthamahaṁ sa ca mama priyaḥ || 7.17 || }
\cquote{Among them, the wise devotee—ever steadfast and single-minded in devotion to Me—is the most distinguished. To them I am most beloved, and they are most beloved to Me.}

\slcol{\Index{उदाराः सर्व एवैते} 
ज्ञानी त्वात्मैव मे मतम् । \\
आस्थितः स हि युक्तात्मा 
मामेवानुत्तमां गतिम् ॥ ७.१८ ॥ }
\translate{udārāḥ sarva ēvaitē 
jñānī tvātmaiva mē matam | 
āsthitaḥ sa hi yuktātmā 
māmēvānuttamāṁ gatim || 7.18 || }
\cquote{Indeed, all these seekers are noble. Yet the wise are regarded as My very Self. Such a person, steadfast in mind, is devoted solely to the supreme goal that is in Me.}

\newpage
\begin{mananam}{\mananamfont Mananam Śloka - 15}
\normalfont\mananamtext What has inspired me to pursue the higher truth? What life experiences have guided me towards seeking God? Can I find gratitude in my past suffering or loss, or in moments of grace or guidance from revered ones, for leading me to the spiritual path? How can I maintain the wisdom to avoid falling back into a state of delusion where I neglect seeking the truth or honouring the Divine?
\end{mananam}
\sreflect
\begin{inspiration}{\mananamfont Inspiration}
\normalfont\mananamtext Many people in this world are content without ever seeking spiritual practices or the truth of a higher being. These individuals, known as \emph{‘pāmara’}, are driven solely by sensual desires. They lack insight into higher realities and have no inclination to seek them. Their focus is entirely on finding happiness in the material world. However, deep within their hearts, there resides an unconscious fear of death and the afterlife, which continually nags at them.
\end{inspiration}
\newpage

\begin{mananam}{\mananamfont Mananam Śloka - 16}
\normalfont\mananamtext Where am I in my journey towards God or Truth? What is my main motivation? Do I only pray to a higher power when I am in distress or need? Do I seek blessings from the great ones for worldly success and protection? Do I have a strong desire to know the Divine or the utmost essence of my own Being? Do I feel reverence or devotion towards those who impart higher wisdom? How can I purify my devotion to God and deepen my understanding of the ultimate Truth?
\end{mananam}
\sreflect
\begin{inspiration}{\mananamfont Inspiration}
\normalfont\mananamtext Even among those who seek God, motivations can differ greatly. Initially, many seekers are driven by the desire to alleviate some form of suffering and attain worldly happiness. Some undertake spiritual practices to gain a competitive edge over others. Eventually, a genuine interest in understanding oneself and the greater force that sustains all creation may develop. With proper guidance, such seekers acquires the right wisdom, which further refines their motivation on the spiritual path. Our goals and desires in life are all based on our level of understanding. 
\end{inspiration}
\newpage


\slcol{\Index{बहूनां जन्मनामन्ते} ज्ञानवान्मां प्रपद्यते । \\
वासुदेवः सर्वमिति स महात्मा सुदुर्लभः ॥ ७.१९ ॥ }
\translate{bahūnāṁ janmanāmantē jñānavānmāṁ prapadyatē | \\
vāsudēvaḥ sarvamiti sa mahātmā sudurlabhaḥ || 7.19 || }
\cquote{After many births, the wise one takes refuge in Me. Realising that \emph{Vāsudeva} is all, such a great soul is exceedingly rare.}

\slcol{\Index{कामैस्तैस्तैर्हृतज्ञानाः} प्रपद्यन्तेऽन्यदेवताः । \\
तं तं नियममास्थाय प्रकृत्या नियताः स्वया ॥ ७.२० ॥ }
\translate{kāmaistaistairhr̥tajñānāḥ prapadyantē:'nyadēvatāḥ | \\
taṁ taṁ niyamamāsthāya prakr̥tyā niyatāḥ svayā || 7.20 || }
\cquote{Driven by one or other desires, people turn to various deities. They follow one or the other practices, as dictated by their tendencies.}

\slcol{\Index{यो यो यां यां} तनुं भक्तः श्रद्धयार्चितुमिच्छति । \\
तस्य तस्याचलां श्रद्धां तामेव विदधाम्यहम् ॥ ७.२१ ॥ }
\translate{yō yō yāṁ yāṁ tanuṁ bhaktaḥ śraddhayārcitumicchati | \\
tasya tasyācalāṁ śraddhāṁ tāmēva vidadhāmyaham || 7.21 || }
\cquote{Whatever form of deity devotees seek to worship with faith, I make that faith steady and unwavering in them.}

\slcol{\Index{स तया श्रद्धया} युक्तस्तस्या राधनमीहते । \\
लभते च ततः कामान्मयैव विहितान्हि तान् ॥ ७.२२ ॥ }
\translate{sa tayā śraddhayā yuktastasyā rādhanamīhatē | \\
labhatē ca tataḥ kāmānmayaiva vihitānhi tān || 7.22 || }
\cquote{Endowed with that faith, the devotees seek the deity of their choice. From that deity they obtain their desired objects, though it is I alone who grant them those fruits.}

\newpage
\begin{mananam}{\mananamfont {Mananam Śloka 18, 19}}
\normalfont\mananamtext How can I enhance my wisdom and gain realisation on the spiritual path? Why is it important to maintain balance between devotion and concentration while also advancing in knowledge? What makes the path of knowledge (\emph{jñāna-yoga}) unique? Why does the Lord extol those on this path to be the dearest and closest to Divine? What is the vision and outlook of a spiritually wise person toward which I can aspire?
\end{mananam}
\sreflect
\begin{inspiration}{\mananamfont Inspiration}
\normalfont\mananamtext The pinnacle of wisdom is realising oneness with the Divine and perceiving that Divinity in everything. Only then can all seeking come to an end. All other forms of worldly wisdom and the forces that sustain it will still uphold a dualistic vision. Where there is duality, the suffering of \emph{saṁsāra} will persist. Such wisdom is attained only when one refuses to compromise with lesser truths and their rewards.
\end{inspiration}
\newpage

\begin{mananam}{\mananamfont Mananam Śloka - 21}
\normalfont\mananamtext What is the nature of my faith? Is it merely a feeling, a sense of loyalty, or blind faith? How did I decide to follow this path? Have I reflected deeply enough on it? Do I observe any progress or positive changes in myself and in my daily life?
Does my intellectual understanding align with my intuitive sense? Does this path offer a vision of the supreme truth and the means to attain it? Or, does it require me to become one-sided and look down on other paths and means?
\end{mananam}
\sreflect
\begin{inspiration}{\mananamfont Inspiration}
\normalfont\mananamtext As one’s understanding deepens, so do one’s aspirations, both in worldly and spiritual life. With deep commitment to any path, our faith intensifies. Consequently, there are various \emph{‘darśanas’} (views or philosophies) even among those who follow the same scriptures, each claiming to hold the ultimate truth. This verse cautions us against placing our faith in the wrong people and means without proper study and discernment. While being faithful to one’s path is commendable, it is problematic if the path is rigid and dismissive of others. In such cases, we must reconsider our approach.
\end{inspiration}
\newpage

\slcol{\Index{अन्तवत्तु फलं तेषां} 
तद्भवत्यल्पमेधसाम् । \\
देवान्देवयजो यान्ति 
मद्भक्ता यान्ति मामपि ॥ ७.२३ ॥ }
\translate{antavattu phalaṁ tēṣāṁ 
tadbhavatyalpamēdhasām | \\
dēvāndēvayajō yānti 
madbhaktā yānti māmapi || 7.23 || }
\cquote{But the results gained by such people of limited understanding is fleeting. Those who worship other deities go to those deities. Yet those who worship Me come (directly) to Me alone.}


\slcol{\Index{अव्यक्तं व्यक्तिमापन्नं} मन्यन्ते मामबुद्धयः । \\
परं भावमजानन्तो ममाव्ययमनुत्तमम् ॥ ७.२४ ॥ }
\translate{avyaktaṁ vyaktimāpannaṁ manyantē māmabuddhayaḥ | 
paraṁ bhāvamajānantō mamāvyayamanuttamam || 7.24 || }
\cquote{Those without wisdom think of Me, the unmanifest, as having assumed manifestation. They do not know My supreme, imperishable, and unsurpassed nature.}

\slcol{\Index{नाहं प्रकाशः सर्वस्य} योगमायासमावृतः । \\
मूढोऽयं नाभिजानाति लोको मामजमव्ययम् ॥ ७.२५ ॥ }
\translate{nāhaṁ prakāśaḥ sarvasya yōgamāyāsamāvr̥taḥ | \\
mūḍhō:'yaṁ nābhijānāti lōkō māmajamavyayam || 7.25 || }
\cquote{I am not manifest to all, being veiled by My divine illusive power (\emph{yoga-māyā}). This deluded world does not know Me, the unborn and imperishable.}

\slcol{\Index{वेदाहं समतीतानि} वर्तमानानि चार्जुन । \\
भविष्याणि च भूतानि मां तु वेद न कश्चन ॥ ७.२६ ॥ }
\translate{vēdāhaṁ samatītāni vartamānāni cārjuna | \\
bhaviṣyāṇi ca bhūtāni māṁ tu vēda na kaścana || 7.26 || }
\cquote{O Arjuna, I know all beings—past, present, and future. Yet none truly knows Me (in My entirety).}

\newpage
\begin{mananam}{\mananamfont Mananam Śloka - 23}
\normalfont\mananamtext To which forms of gods or deities do I pray? What outcomes or rewards do I seek from these practices? Are the promised results of my practices finite and time-bound? Is my devotion directed towards lesser gods or to the Supreme? Do I have a sufficient understanding of the Supreme Being, the Self, and the ultimate Truth?
\end{mananam}
\sreflect
\begin{inspiration}{\mananamfont Inspiration}
\normalfont\mananamtext Regardless of the nature of our practices, we should aspire to realize the Supreme Being, which is the ultimate reality of our own existence. All other forms of worship and practices will only provide short-term rewards. True wisdom, acquired through a lineage of \emph{gurus} and the scriptures, ensures that we remain on the right path and that our efforts will not be diverted from the supreme goal of life.
\end{inspiration}
\newpage

\slcol{\Index{इच्छाद्वेषसमुत्थेन} द्वन्द्वमोहेन भारत । \\
सर्वभूतानि संमोहं सर्गे यान्ति परन्तप ॥ ७.२७ ॥ }
\translate{icchādvēṣasamutthēna dvandvamōhēna bhārata | \\
sarvabhūtāni saṁmōhaṁ sargē yānti parantapa || 7.27 || }
\cquote{O Bhārata, O Parantapa (Arjuna), all beings are deluded at birth by the dualities of desire and aversion, arising from delusion.}

\slcol{\Index{येषां त्वन्तगतं पापं} जनानां पुण्यकर्मणाम् । \\
ते द्वन्द्वमोहनिर्मुक्ता 
भजन्ते मां दृढव्रताः ॥ ७.२८ ॥ }
\translate{yēṣāṁ tvantagataṁ pāpaṁ 
janānāṁ puṇyakarmaṇām | \\
tē dvandvamōhanirmuktā 
bhajantē māṁ dr̥ḍhavratāḥ ||7.28 || }
\cquote{Freed from sin, purified by righteous action, and freed from the delusion of dualities, such steadfast ones worship Me with steadfast vows.}

\slcol{\Index{जरामरणमोक्षाय} मामाश्रित्य यतन्ति ये । \\
ते ब्रह्म तद्विदुः कृत्स्नमध्यात्मं कर्म चाखिलम् ॥ ७.२९ ॥ }
\translate{jarāmaraṇamōkṣāya māmāśritya yatanti yē | \\
tē brahma tadviduḥ kr̥tsnamadhyātmaṁ karma cākhilam || 7.29 || }
\cquote{Those who strive for liberation from old age and death, taking refuge in Me, they know \emph{Brahman}, the Self, and all action in its entirety.}

\slcol{\Index{साधिभूताधिदैवं मां} साधियज्ञं च ये विदुः । \\
प्रयाणकालेऽपि च मां ते विदुर्युक्तचेतसः ॥ ७.३० ॥ }
\translate{sādhibhūtādhidaivaṁ māṁ sādhiyajñaṁ ca yē viduḥ | \\
prayāṇakālē:'pi ca māṁ tē viduryuktacētasaḥ || 7.30 || }
\cquote{Those who know Me along with the knowledge of the material elements (\emph{ādhibhūta}), the divine beings (\emph{ādhidaiva}), and sacrifice (\emph{ādhiyajña})—such steadfast ones know Me even at the time of death.}

\newpage
\begin{mananam}{\mananamfont Mananam Śloka - 29}
\normalfont\mananamtext Who and what are my sources of support in the world? In which areas can I rely on my family, friends, and others for support? In which areas might I not be able to rely on them? Under what conditions are worldly sources of refuge inadequate or limited? What concept can I envision for a refuge that transcends these limitations? How can such an Infinite Being be the ultimate assurance, enhancing my sense of self?
\end{mananam}
\sreflect 
\begin{inspiration}{\mananamfont Inspiration}
\normalfont\mananamtext In many aspects of material life, we can rely on support from people, society, and the government. However, when facing the inevitability of old age and death, our highest refuge is God alone. Regardless of how much we rationalise, the mind needs something to hold onto. For a practitioner, the form of God or a felt presence of the Divine provides that support. For one who is stabilised in wisdom, the realisation of the Self and its unbroken connection with the Infinite serves as that anchor.
\end{inspiration}
\newpage

\begin{center}
\devfont ॐ तत्सदिति श्रीमद्भगवद्गीतासूपनिषत्सु  ब्रह्मविद्यायां योगशास्त्रे\\ श्रीकृष्णर्जुनसंवादे ज्ञानविज्ञानयोगो  नाम सप्तमोऽध्यायः॥\\
\chapEndSlokaTranslate{ōṃ tatsaditi śrīmadbhagavadgītās ūpaniṣatsu \\brahmavidyāyāṃ yōgaśāstrē\\ śrīkṛṣṇārjunasaṃvādē jñānavijñānayōgō\\nāma prathamo'dhyāyaḥ ‖}
\end{center}